\documentclass{article}
\usepackage[T1]{fontenc}
\usepackage{lmodern}
\usepackage{graphicx}
\usepackage{amsmath,amssymb}
\usepackage[a4paper,margin=2cm]{geometry}
\usepackage[absolute,overlay]{textpos}
\setlength{\TPHorizModule}{1mm}
\setlength{\TPVertModule}{1mm}
\usepackage[indonesian]{babel}
\usepackage{hyperref}
\setlength{\emergencystretch}{2em}
% unit: Halaman 57

\begin{document}

% === Halaman 57 (python/native) ===

57

• Contoh: 

      1 


            2 


      3 


            4 

        5 

 LJVBQSTNEZLQMEDLJVMAMPKAUFAVATLJVDAYYVNFJQLNPLJVHKVTRNFLJVCMLKETA 

LJVHUYJVSFKRFTTWEFUXVHZNP

  6 



 Kriptogram yang berulang adalah LJV   


Jarak LJV ke-1 dengan LJV ke-2 = 15 



Jarak LJV ke-2 dengan LJV ke-3 = 15 


Jarak LJV ke-3 dengan LJV ke-4 = 15


Jarak LJV ke-4 dengan LJV ke-5 = 10


Jarak LJV ke-5 dengan LJV ke-6 = 10

 Faktor pembagi 15 = \{3, 5, 15\}

 Faktor pembagi 10 = \{2, 5, 10\}

 Irisan kedua himpunan ini = 5. Jadi, panjang kunci diperkirakan = 5

\end{document}
